\documentclass[11pt]{article}
\usepackage{fontspec}
\usepackage[french]{babel}
\usepackage[leqno]{amsmath}
\usepackage{amssymb}
\usepackage{amsthm}
\usepackage{mathrsfs}
\usepackage{tikz-cd}
\usepackage{enumitem}

\frenchsetup{AutoSpacePunctuation=false}

\newcommand{\printedpage}[1]{\marginpar{\footnotesize p.~#1}}
\newcommand{\msp}[1]{\marginpar{\footnotesize [#1]}}
\newcommand{\cM}{\mathcal{M}}
\newcommand{\cF}{\mathcal{F}}
\newcommand{\Ho}{\mathrm{Ho}}
\newcommand{\ho}{\mathrm{ho}}
\newcommand{\Fl}{\mathrm{Fl}}
\newcommand{\Fib}{\underline{\mathrm{Fib}}}
\newcommand{\Wu}{W^{\mathrm{u}}}
\newcommand{\lc}{\mathrm{lc}}
\newcommand{\circled}[1]{\tikz[baseline=(c.base)]\node[draw,circle,inner sep=1pt](c){#1};}

\begin{document}

\printedpage{135}
% head: A. Grothendieck : Dérivateurs, Ch. XII, Édition des Universités de Montpellier II et Paris VII

\begin{quote}
transforme $W_X$ dans $W_Y$, où $W_X \subseteq \Fl(\cM/X)$ \texttt{[est]} défini par

$u \in W_X \iff \sigma(U) \in W$ \ ($\sigma : \cM/X \longrightarrow \cM$ le foncteur ‘source’, ou ‘oubli de $X$’). % [sic] \sigma(U) -> \sigma(u), the arrow is called u
\end{quote}

(\footnote[120]{\textbf{NB} Donner ici la proposition: $\mathrm{Fib}_W$ \textit{stable par composition et par changement de base, et contient les isomorphismes, stable par facteurs directs.}}). Posons, pour tout $S$ dans $\cM$,
\[
\Ho^!_W(S) = (\cM/S)W_S^{-1},
\]
qui joue le rôle d'une catégorie de $W$-types d'homotopie relatifs \textit{localement constants} sur $S$. On aimerait décrire une loi contravariante en $S$ pour $\Ho^!_W(S)$, i.e. pour
\[
f : S' \longrightarrow S
\]
flèche quelconque dans $\cM$, définir
\[
f^*_W : \Ho^!_W(S) \xrightarrow{?} \Ho^!_W(S'). \tag*{$(*)$}
\]
La difficulté est, bien sûr, que le foncteur changement de base
\[
f^* : \cM/S \longrightarrow \cM/S'
\]
n'est \textit{pas} en général compatible avec les \msp{page 209} localiseurs $W_S$, $W_{S'}$. La notion de $W$-fibration sert donc en premier lieu à dégager la condition la plus évidente sur une flèche $f : S' \longrightarrow S$ dans $\cM$, pour que le foncteur changement de base soit compatible avec les localiseurs $W_S$, $W_{S'}$ (c'est \textit{ça} la définition), donc induit $(*)$ sur les localisées.

Le \textit{‘problème des fibres homotopiques’} est justement de définir $(*)$ pour \textit{toute} flèche dans $\cM$, pas seulement pour les $W$-fibrations, de tellle façon que les $\Ho^!_W(S)$ forment, pour $S$ % [sic] tellle -> telle, letter doubled
variable, les fibres d'une catégorie fibrée remarquable sur $\cM$. Quand $\cM$ admet un objet final, de sorte que l'on peut regarder les flèches
\[
p : e \longrightarrow S \qquad (p \text{ comme ‘point’})
\]
dans $\cM$ comme des ‘points’ de $S$, la ‘fibre $W$-homotopique’ d'un objet quelconque $X/S$ de $\cM/S$ en ledit point sera par définition l'élément de
\[
\Ho_W = \Ho_W(\cM) \simeq \Ho^!_W(e)
\]
défini par
\[
p^*(\ho^!_{W/S}(X/S)),
\]
\printedpage{136}
% head: A. Grothendieck : Dérivateurs, Ch. XII, Édition des Universités de Montpellier II et Paris VII
\msp{page 210}
où
\[
\ho^!_{W/S} : \cM/S \longrightarrow \Ho^!_W(S) \overset{\text{déf}}{=} (\cM/S)(W_S)^{-1}
\]
est le foncteur de localisation canonique.

Pour tout $S$ dans $\cM$, soit
\[
\Fib_W S \subseteq \cM/S
\]
la sous-catégorie pleine formée des objets $X/S$ de $\cM/S$ tels que $X \longrightarrow S$ soit $W$-fibrant. On va définir une variante
\[
\Ho^{\lc}_W(S)
\]
de $\Ho^!_W(S)$, comme
\[
\Ho^{\lc}_W(S) \overset{\text{déf}}{=} (\Fib_W S)(\Wu_S)^{-1},
\]
où

\begin{quote}
$\Wu_S \subseteq \Fl(\Fib_W S)$ est le localiseur formé des flèches $u : X \longrightarrow Y$ qui sont ‘universellement sur $S$ dans $W$’, i.e. telles que pour toute flèche $S' \longrightarrow S$ de $\cM$ de but $S$, la flèche déduite par changement de base,
\[
u' : X' = X \times_S S' \longrightarrow Y' = Y \times_S S'
\]
est dans $W_{S'}$.
\end{quote}

On a ‘fait ce qu'il fallait’ pour que $\Ho^{\lc}_W(S)$ soit contravariant en $S$ pour des flèches \textit{quelconques} $f : S' \longrightarrow S$ de $\cM$, elles définissent
\[
f^{*\,\lc}_W : \Ho^{\lc}_W(S) \longrightarrow \Ho^{\lc}_W(S'), \tag*{$(**)$}
\]
\msp{page 211}
puisque les $W$-fibrations sont stables par changement de base, et que par suite $f^* : \cM/S \longrightarrow \cM/S'$ induit
\[
f^*_{\Fib_W} : \Fib_W S \longrightarrow \Fib_W S'
\]
qui, par définition de $\Wu_S$ et $\Wu_{S'}$, est compatible avec les localiseurs
\[
f^*_{\Fib_W} : (\Fib_W S, \Wu_S) \longrightarrow (\Fib_W S', \Wu_{S'}),
\]
et passe donc aux catégories des fractions en une flèche $(*)$.

On peut se demander pourquoi ne pas être plus généreux, et définir $\Ho^{\lc}_W(S)$ comme $(\cM/S)(\Wu_S)^{-1}$, en définissant $\Wu_S$ dans $\cM/S$ de la même façon que tantôt dans $\Fib_W S$ - ce qui ne coûte pas plus cher. Une première réponse est d'expérience le localiseur $\Wu_S$ dans $\cM/S$ tout entier s'avère pratiquement inutilisable, d'une part parce que l'appartenance à $\Wu_S$ (pour une flèche entre objets de $\cM/S$ qui ne sont pas très particuliers, p.ex. $W$-fibrés sur $S$) est invérifiable (sauf pour des isomorphismes et des choses comme ça \dots), et que d'autre part la catégorie de fractions est prohibitivement
\printedpage{137}
% head: A. Grothendieck : Dérivateurs, Ch. XII, Édition des Universités de Montpellier II et Paris VII
\msp{page 212}
grosse et ne peut servir (semble-t-il) à rien. Une autre réponse, c'est qu'on veut que la catégorie obtenue réponde bien à l'intuition de types d'homotopie relatifs \textit{localement constants} sur $S$ - et on veut donc que pour un changement de base quelconque $f : S' \longrightarrow S$ par un

\begin{quote}
$S'$ dans $\cM/S$, d'où un foncteur composé
\[
\begin{tikzcd}[column sep=small]
\cF_S(\Wu_S)^{-1} \arrow[r, "f^{*\,\lc}_W"] \arrow[rr, dotted, bend right=15, "\varphi_{S'}"'] & \cF_{S'}(\Wu_{S'})^{-1} \arrow[r] & \Ho_W
\end{tikzcd}
\]

(\footnote[121]{$\cF$ désigne ici une sous-catégorie fibrée de $\cM/\cM = \underline{\Fl}(\cM)$ sur $\cM$, substitut éventuel de celle formée de $\Fib_W S$ ou de $\cM/S$ tout entier.}) déduit des foncteurs de catégories avec localiseurs
\[
(\cF_S, \Wu_S) \xrightarrow{f^*} (\cF_{S'}, \Wu_{S'}) \xrightarrow{\text{oubli de } S'} (\cM, W)
\]
\end{quote}

(foncteur composé $\varphi_{S'}$ qui peut être vu comme le foncteur \textit{‘fibre homotopique en l'objet $S'$ sur $S$’}, cet objet étant visualisé comme un ‘point’ de $S$ ‘à valeurs dans $S'$ ’ suivant le yoga fonctoriel général \dots), le foncteur
\[
\begin{aligned}
S' &\longmapsto \varphi_{S'}(\xi) \qquad \text{pour } \xi \text{ dans } \cF_S(\Wu_S)^{-1}\\
\cM/S &\longrightarrow \Ho_W
\end{aligned}
\]
transforme flèches dans $W_S$ en isomorphismes : la ‘fibre $W$-homotopique’ en $S'/S$ d'un objet de $\cF_S(\Wu_S)^{-1}$ (vu comme $W$-type d'homotopie \textit{localement constant} sur $S$) doit être ‘la même’, si $S'$ remplacé par $S''$ avec $S'' \longrightarrow S'$ donné et dans $W_S$. C'est pour cela qu'il est raisonnable de prendre $\cF_S$ ‘assez petit’ (savoir $\Fib_W S$, qui est \msp{page 213} exactement ajusté pour cela), plutôt que de prendre $\cM/S$ tout entier.

On peut donc dire que la catégorie $\Fib_W S$, et le localiseur $\Wu_S$ ont été construits ad-hoc, le plus simplement possible, pour donner lieu à un localisé $\Ho^{\lc}_W(S)$ qui réponde à l'intuition de $W$-types d'homotopie relatifs localement constants des deux façons suivantes :

\begin{enumerate}[label={}]
\item[1°)] Il y a spontanément une loi contravariante en $S$, pour des flèches quelconques $S' \longrightarrow S$ dans Cat (on a défini $\Wu_S$ juste exprès en conséquence).
\item[2°)] que la notion de ‘fibre $W$-homotopique en $S'/S$’ d'un objet $\xi$ de $\Ho^{\lc}_W$ satisfasse à la condition précédente, i.e. une flèche $S'' \longrightarrow S'$ dans $\cM/S$ qui est dans $W_S$ définisse un \textit{isomorphisme} (dans $\Ho_W$) des fibres homotopiques
\[
\underbrace{\xi_{S''}}_{\varphi_{S''}(\xi)} \longrightarrow \underbrace{\xi_{S'}}_{\varphi_{S'}(\xi)}
\]
(et on a choisi $\cF_S = \Fib_W(S)$ juste exprès pour ça).
\end{enumerate}

\printedpage{138}
% head: A. Grothendieck : Dérivateurs, Ch. XII, Édition des Universités de Montpellier II et Paris VII

On a donc une excellente notion de fibres $W$-homotopiques pour des $X$ $W$-fibrés sur $S$ - cette fois excellente pour des changements de base \textit{quelconque} - et il n'y a aucun mérite à ça puisqu'on a tout défini exprès pour ça! % [sic] quelconque -> quelconques, agreement with "des changements"

\msp{page 214}
Pour le dire autrement, on a défini \textit{deux} notions de $W$-types d'homotopie relatifs sur $S$, l'une de façon à donner un tel type d'homotopie $\ho^!_{W,S}(X)$ pour \textit{tout} objet de $\cM/S$, mais il n'y a a priori une bonne notion de changement de base que pour les changements de base $S' \longrightarrow S$ qui sont $W$-fibrés, l'autre, qui correspond plus directement à l'intuition d'un $W$-type d'homotopie localement constant sur $S$, n'associe un tel type $\ho^{\lc}_{W,S}(X)$ qu'aux $X$ de $\cM/S$ qui sont $W$-fibrés sur $S$, en revanche cette notion admet une bonne notion d'image inverse et de ‘fibre $W$-homotopique’ pour des changements de base quelconques.

Une ‘bonne’ théorie des fibres $W$-homotopiques doit réunir en une seule les deux notions précédentes, de sorte que les avantages des deux s'ajoutent, et les désavantages s'éliminent! Elle s'obtient à l'aide d'un axiome draconien sur la situation (alors que jusqu'à maintenant on n'a rien utilisé sur $W \subseteq \Fl(\cM)$, si ce n'est tout au plus la saturation faible, en plus de l'existence des produits fibrés) :

\msp{page 215}
\circled{a} (\textit{Axiome de la fibre homotopique.}) Pour tout objet $S$ dans $\cM$, le foncteur canonique
\[
\Ho^{\lc}_W(S) \longrightarrow \Ho^!_W(S)
\]
déduit de l'inclusion
\[
(\Fib_W(S), \Wu_S) \longrightarrow (\cM/S, W_S)
\]
est une \textit{équivalence de catégories.}

Cet axiome implique que dans $\Fib_W S$, $\Wu_S$ et le localiseur induit par $W_S$ (que je désigne encore par $W_S$ par abus de notation) ont même saturé fort. En fait, il est naturel de compléter l'axiome précédent par le suivant.

\circled{b} (Axiome adjoint de l'axiome de la fibre homotopique.) Sur $\Fib_W S$, les localiseurs $\Wu_S$ et $W_S$ coïncident, i.e. on a ceci :

\begin{quote}
Soient $X$, $Y$ dans $\Fib_W S$, $X \xrightarrow{u} Y$ un $S$-morphisme tel que $u \in W_S$, alors $u \in \Wu_S$ (i.e. le morphisme $u' : X' \longrightarrow Y'$ déduit par tout changement de base $S' \longrightarrow S$ est encore dans $W$).
\end{quote}

\msp{page 216}
Je présume que a) sera rarement satisfait sans que b) le soit (et inversement d'ailleurs), et que la démonstration de a) et b) se fera pratiquement en une seule haleine.

Si $\cM = \mathrm{Cat}$, $W$ un localiseur fondamental, alors la notion adaptée ici de $W$-fibration (lire : ‘ultraforte’, dans la terminologie spécifique développée dans Cat) est à tel point forte, qu'elle implique déjà pratiquement b) (moyennant W(4a$'$), il est vrai, cf. p. 53), et
\printedpage{139}
% head: A. Grothendieck : Dérivateurs, Ch. XII, Édition des Universités de Montpellier II et Paris VII
a) n'est guère vérifiable que si on a W(4a) (\footnote[122]{Non, il me semble que W(4b) doit suffire \dots}) (loc. cit.). Donc ici c'est ‘l'axiome des fibrations’ sur $W$, sous sa forme forte W(4a), qui est la clef de voûte d'une bonne théorie des fibres homotopiques.

Une telle théorie (revenant au cas général) revient à dire, pour l'essentiel, qu'il y a ‘assez’ de $W$-fibrations (‘ultrafortes’) pour pouvoir décrire \textit{tout} objet d'un $\cM/S$, à $W_S$-équivalence près, comme un objet $W$-fibrant sur $S$. Dans les cas à ma connaissance, on a de plus ceci : on \msp{page 217} a pour tout $S$ dans $\cM$ une sous-catégorie $\cF_{0\,S}$ de $\Fib_W S$ (je pense par exemple à $\underline{\mathrm{Parf}}_W S$, voire $\underline{\mathrm{Parf}}_{W_\omega} S$, dans le cas où $\cM = \mathrm{Cat}$), telle que pour toute sous-catégorie pleine $\cF$ de $\cM/S$ telle que
\[
\cF_{0\,S} \subseteq \cF \subseteq \cM/S
\]
(a fortiori, pour une $\cF$ pleine telle que $\Fib_W S \subseteq \cF \subseteq \cM/S$), les inclusions induisent des équivalences de catégories
\[
\cF_{0\,S}(W_S)^{-1} \xrightarrow{\sim} \cF(W_S)^{-1} \xrightarrow{\sim} \underbrace{(\cM/S)(W_S)^{-1}}_{\overset{\text{déf}}{=} \mathrm{Hot}^!_W(S)}. % [sic] Hot -> Ho, elsewhere the notation is Ho^!_W(S)
\]
Le choix de $\cF_{0\,S}$ ne semble pas pouvoir se fixer de façon ‘la meilleure possible’, une fois pour toutes dans le contexte axiomatique adopté ici. C'est un élément de structure supplémentaire dans $\cM$ (en plus du localiseur $W$), au même titre que la donnée des classes de $Q$-fibrations et $Q$-cofibrations, à la \textsc{Quillen}. J'ignore si les axiomes a), b) ci-dessus (axiome de la fibre homotopique) impliquent déjà \msp{page 218} la variante de \circled{a} :

\circled{c} Pour toute sous-catégorie pleine $\cF$ de $\cM/S$ telle que
\[
\Fib_W S \subseteq \cF \subseteq \cM/S
\]
les inclusions induisent des équivalences
\[
\underbrace{(\Fib_W S)(W_S)^{-1}}_{\begin{subarray}{l} = \ \Ho^{\lc}_W(S)\\ \text{moyennant a)}\end{subarray}} \xrightarrow{\sim} \cF(W_S)^{-1} \xrightarrow{\sim} \underbrace{(\cM/S)(W_S)^{-1}}_{=\Ho^!_W(S)}.
\]

Mais je présume que pratiquement, quand on pourra vérifier \circled{a}, on pourra vérifier \circled{c} également (\footnote[123]{Cf. page 240 ff. Il faudrait ici développer un peu plus le formalisme de $\Ho^{\lc}_W(S)$ et les fibres homotopiques, sous l'hypothèse \circled{a}\,\circled{b} \ plus haut.}).

\end{document}
